% hierarchical_coupling.tex
% Attention as Renormalization
\documentclass[11pt,a4paper]{article}

% ── Packages ──────────────────────────────────────────────────
\usepackage[utf8]{inputenc}
\usepackage{amsmath,amssymb,amsthm}
\usepackage{graphicx}
\usepackage{booktabs}
\usepackage{multirow}
\usepackage{xcolor}
\usepackage[margin=2.5cm]{geometry}
\usepackage{hyperref}
\usepackage{natbib}
\usepackage{cleveref}
\usepackage{algorithm}
\usepackage{algpseudocode}
\usepackage{caption}
\usepackage{subcaption}

% ── Notation ──────────────────────────────────────────────────
\newcommand{\field}{\phi}
\newcommand{\vfield}{\boldsymbol{\phi}}
\newcommand{\deposit}{d}
\newcommand{\vdeposit}{\mathbf{d}}
\newcommand{\evap}{\varepsilon}
\newcommand{\retain}{\alpha}
\newcommand{\R}{\mathbb{R}}
\newcommand{\E}{\mathbb{E}}
\newcommand{\x}{\mathbf{x}}
\newcommand{\q}{\mathbf{q}}
\newcommand{\K}{\mathbf{k}}
\newcommand{\vv}{\mathbf{v}}
\newcommand{\bigo}{\mathcal{O}}

\newtheorem{hypothesis}{Hypothesis}
\newtheorem{proposition}{Proposition}
\newtheorem{definition}{Definition}
\newtheorem{remark}{Remark}

% ── Title ─────────────────────────────────────────────────────
\title{%
  Attention as Renormalization:\\
  Multi-Scale Vector Fields for\\
  Length-Invariant Sequence Processing%
}

\author{
  Draft --- Work in Progress
}

\date{\today}

\begin{document}
\maketitle

% ══════════════════════════════════════════════════════════════
\begin{abstract}
Length generalization---the ability to process sequences longer than
those seen in training---is a scale-invariance problem.  We argue
that multi-scale recurrent fields coupled across transformer layers
implement a discrete renormalization group (RG) hierarchy, where
coarse-graining over successively longer timescales produces
representations that are invariant to total sequence length.  We
develop this idea in two stages.  First, we show that \emph{scalar}
pheromone fields biasing attention logits exhibit an all-or-nothing
ablation structure: five of six partial configurations fail, and only
the complete system (feedback deposits, fast/slow fields, cross-layer
coupling) generalizes---but poorly compared to dedicated recurrent
models.  Second, we upgrade from scalar logit biases to \emph{vector
social fields} that shift attention keys, drawing on a companion
dual-channel architecture.  The vector field provides dramatically
higher state capacity (cumsum mod~16 at $2\times$: 0.999 vs 0.279
baseline) and a natural geometric interpretation: the field modulates
\emph{what attention looks for}, not just \emph{how much} it attends.
The renormalization framing unifies these results: the fast field
captures token-level dynamics, the slow field coarse-grains over
temporal blocks, and cross-layer coupling iterates the RG
transformation.  Length generalization emerges when the system
operates near a critical fixed point where these scales interact
self-consistently.
\end{abstract}


% ══════════════════════════════════════════════════════════════
\section{Introduction}
\label{sec:intro}

Transformers do not generalize across sequence lengths.  This is not
merely empirical: bounded-depth transformers provably cannot simulate
automata or track state over arbitrary-length inputs
\citep{merrill2023transformers}.  The core issue is that attention
patterns are tied to positional representations that do not
extrapolate.

We reframe this as a \textbf{scale-invariance problem}.  A model that
generalizes to longer sequences has learned representations that look
the same regardless of total length---just as a system at a critical
point looks the same regardless of magnification.  Length
generalization is scale invariance in time.

This framing suggests a constructive approach: \emph{build a
renormalization hierarchy into the attention mechanism}.  In
condensed-matter physics, the renormalization group (RG) achieves
scale invariance by iteratively coarse-graining short-distance
structure while preserving long-distance behavior
\citep{kadanoff1966,wilson1971rg}.  The key insight is that at a
critical fixed point, the system is self-similar across scales---the
coarse-grained version is statistically identical to the original.

We implement a discrete RG hierarchy within transformer attention
using \emph{stigmergic fields}---vector-valued recurrent states
inspired by ant pheromone trails \citep{antpaper}.  The hierarchy has
three levels:

\begin{enumerate}
  \item \textbf{Fast field} ($\tau_f \approx 6$ tokens): captures
    local, token-level dynamics.  This is the ``microscopic'' scale.
  \item \textbf{Slow field} ($\tau_s \approx 50$ tokens):
    coarse-grains the fast field over temporal blocks.  This is a
    \emph{block-spin transformation in time}---the temporal analogue
    of Kadanoff's block-spin RG \citep{kadanoff1966}.
  \item \textbf{Cross-layer coupling}: the slow field's terminal
    state seeds the next layer's initial condition.  Each layer
    operates at a successively coarser temporal scale.  This is
    \emph{iteration of the RG transformation}.
\end{enumerate}

The field enters attention not as a scalar logit bias (our earlier
approach, which we show is too weak) but as an \emph{additive key
shift}---a vector that modulates the geometry of attention.  This
means the coarse-grained field changes \emph{what the head looks for},
not just \emph{how much it attends}, providing the state capacity
needed for multi-bit sequence tasks.

\paragraph{Contributions.}
\begin{enumerate}
  \item A renormalization-group interpretation of multi-scale temporal
    coupling in attention, connecting length generalization to scale
    invariance at a critical fixed point.
  \item A controlled ablation showing that the RG hierarchy is
    all-or-nothing: five of six partial configurations fail, mapping
    onto the concept of relevant vs.\ irrelevant operators in RG.
  \item Evidence from two complementary architectures: scalar fields
    (clean ablation, limited capacity) and vector fields (high
    capacity, dramatic length generalization on 4-bit tasks).
  \item A structural analysis of why the minimum viable system requires
    feedback, scale separation, and vertical coupling---three
    components that together implement the renormalization hierarchy.
\end{enumerate}


% ══════════════════════════════════════════════════════════════
\section{Background}
\label{sec:background}

\subsection{The Length Generalization Problem}

Bounded-depth transformers with finite precision cannot simulate
general finite automata \citep{merrill2023transformers}.  For tasks
like parity---computing the running sum modulo 2---a transformer
trained on length-$T$ inputs will memorize positional heuristics that
break at length $T' > T$.

Existing approaches include positional extrapolation (ALiBi;
\citealt{press2022alibi}), chain-of-thought, SSMs
(S4~\citealt{gu2022s4}; Mamba~\citealt{gu2023mamba}), and
interleaved hybrids (Jamba~\citealt{lieber2024jamba};
Griffin~\citealt{de2024griffin}).  None directly test what
\emph{structure} of recurrence enables length generalization when
coupled into the attention mechanism itself.

\subsection{Renormalization Group Basics}

The renormalization group \citep{kadanoff1966,wilson1971rg} is a
framework for understanding systems at multiple scales.  The
procedure is:

\begin{enumerate}
  \item \textbf{Coarse-grain}: partition the system into blocks,
    replace each block with a single effective degree of freedom.
  \item \textbf{Rescale}: the block-level system has the same form
    as the original but with renormalized parameters.
  \item \textbf{Iterate}: repeated coarse-graining flows through
    parameter space toward fixed points.
\end{enumerate}

At a \emph{critical fixed point}, the system is scale-invariant: it
looks statistically identical at all magnifications.  The critical
point is characterized by:
\begin{itemize}
  \item \textbf{Relevant operators}: perturbations that grow under
    RG flow.  Removing a relevant operator drives the system away
    from the fixed point.
  \item \textbf{Irrelevant operators}: perturbations that shrink
    under RG flow.  The system self-corrects.
  \item \textbf{Universality}: systems with the same relevant
    operators share the same critical behavior, regardless of
    microscopic details.
\end{itemize}

\subsection{Stigmergy and Ant Trail Formation}

Stigmergy is the indirect coordination mechanism used by social
insects: agents deposit pheromones that influence future behaviour.
\citet{antpaper} identify a three-ingredient recipe for edge-of-chaos
dynamics in lattice ant models: \emph{amplification} (positive
feedback), \emph{injection} (activity-dependent deposition), and
\emph{erosion} (uniform decay).  Removing any ingredient causes
collapse.  Crucially, erosion must be \emph{uncoupled} from the field
it erodes---a structural condition that we will interpret as a
requirement for well-defined coarse-graining.


% ══════════════════════════════════════════════════════════════
\section{Architecture: A Renormalization Hierarchy in Attention}
\label{sec:architecture}

\subsection{Scalar Fields: The Minimal Case}
\label{sec:scalar-field}

We begin with the simplest version to establish the ablation
structure.  For each attention head $h$, define a scalar field
$\field_i^{(h)}$ at position $i$:
\begin{equation}
  \field_i^{(h)} = \retain \cdot \field_{i-1}^{(h)} + \deposit_i^{(h)},
  \qquad \retain = 1 - \evap
  \label{eq:scalar-recurrence}
\end{equation}
The field enters as an additive logit bias:
\begin{equation}
  \ell_{ij}^{(h)}
    = \frac{\q_i^{(h)\top} \K_j^{(h)}}{\sqrt{d_h}}
    + \field_j^{(h)}
  \label{eq:scalar-logit}
\end{equation}

The deposit source determines the field's expressive power.
\emph{Input deposits} ($\deposit_i = W_d \x_i$) make the field a
deterministic moving average---provably inert for state tracking.
\emph{Feedback deposits} ($\deposit_i = W_d \mathbf{o}_i$, where
$\mathbf{o}_i$ is the attention output) close the loop:
\begin{equation}
  \field \xrightarrow{\text{bias}} \text{attention}
  \xrightarrow{\text{aggregate}} \text{output}
  \xrightarrow{W_d} \text{deposit}
  \xrightarrow{\text{accumulate}} \field
  \label{eq:feedback-loop}
\end{equation}
creating a genuine nonlinear recurrence.

\paragraph{Limitation.}  The scalar field says ``attend more here''
but not ``attend in this direction.''  With $H$ heads, the field
provides $O(H)$ bits of recurrent state---enough for 1-bit parity
but insufficient for richer tasks.

\subsection{Vector Fields: Modulating Attention Geometry}
\label{sec:vector-field}

We replace the scalar field with a \emph{vector social field}
$\vfield_i^{(h)} \in \R^{d_f}$ that modulates keys directly:

\begin{definition}[Vector social field]
For each head $h$ with key dimension $d_h$ and field dimension
$d_f = d_h$:
\begin{align}
  \vfield_i^{(h)} &= \retain \cdot \vfield_{i-1}^{(h)}
    + \vdeposit_i^{(h)}
    \label{eq:vector-recurrence} \\
  \vdeposit_i^{(h)} &= W_{\mathrm{dep}}^{(h)} \, \mathbf{o}_i^{(h)}
    \label{eq:vector-deposit}
\end{align}
where $\mathbf{o}_i^{(h)}$ is the attention output at position $i$,
head $h$.
\end{definition}

The field enters attention through additive key modulation:
\begin{equation}
  \K_j'^{(h)} = \K_j^{(h)} + W_{\mathrm{read}}^{(h)} \, \vfield_j^{(h)}
  \label{eq:key-modulation}
\end{equation}

This is a qualitatively different mechanism.  The scalar field biases
\emph{how much} a position is attended to.  The vector field shifts
keys in a learned direction, changing \emph{what kind of match}
occurs.  With field dimension $d_f = d_h$, each head carries
$d_f$-dimensional recurrent state---orders of magnitude more than the
scalar case.

\paragraph{Biological interpretation.}  Ant pheromones are not scalar
intensities---they are chemical blends whose composition encodes
qualitative information (food type, danger, nest direction).  A vector
field is the natural analogue: the field's direction encodes
\emph{what kind} of signal was deposited, while its magnitude encodes
strength.

\subsection{The Renormalization Hierarchy}
\label{sec:rg-hierarchy}

We now build the multi-scale structure, interpreted as a discrete
RG transformation.

\subsubsection{Level 0: Token-Level Dynamics (Input)}

The input sequence $\x_1, \ldots, \x_T$ defines the ``microscopic''
degrees of freedom.  Standard attention operates at this level:
each query $\q_i$ matches against all keys $\K_j$ for $j \leq i$.

\subsubsection{Level 1: Fast Field (Local Coarse-Graining)}

The fast field $\vfield^{\mathrm{f}}$ with evaporation rate
$\evap_f = 0.15$ ($\tau_f \approx 6$) performs the first
coarse-graining.  At each position, it accumulates
attention-derived deposits over a window of $\sim\!\tau_f$ tokens:
\begin{equation}
  \vfield^{\mathrm{f},(L)}_i
    = \retain_f \cdot \vfield^{\mathrm{f},(L)}_{i-1}
    + \vdeposit_i^{(L)}
  \label{eq:fast}
\end{equation}

This is the temporal analogue of averaging over a spatial block in
the Kadanoff construction: short-timescale fluctuations are smoothed
into a local summary vector that modulates keys at nearby positions.

\subsubsection{Level 2: Slow Field (Block-Spin Transformation)}

The slow field $\vfield^{\mathrm{s}}$ with $\evap_s = 0.02$
($\tau_s \approx 50$) coarse-grains the fast field:
\begin{equation}
  \vfield^{\mathrm{s},(L)}_i
    = \retain_s \cdot \vfield^{\mathrm{s},(L)}_{i-1}
    + h^{(L)}\!\left(\vfield^{\mathrm{f},(L)}_i\right)
  \label{eq:slow}
\end{equation}
where $h^{(L)}: \R^{d_f} \to \R^{d_f}$ is a learned linear
projection.

This is the \textbf{block-spin transformation}: the slow field
replaces blocks of $\sim\!\tau_s/\tau_f \approx 8$ fast-field values
with a single effective vector.  The projection $h^{(L)}$ is the
``block-spin rule''---it determines which components of the
fast-field structure survive coarse-graining.

The rescaling factor $b = \tau_s / \tau_f$ controls the RG step
size.  Our values give $b \approx 8$, meaning each slow-field
update integrates over $\sim$8 fast-field correlation lengths.

\subsubsection{Level 3: Cross-Layer Coupling (Iterating the RG)}

The slow field's terminal state at layer $L$ seeds layer $L{+}1$:
\begin{equation}
  \vfield^{\mathrm{s},(L+1)}_0
    = \sigma\!\left(w_g^{(L)}\right) \cdot
      P^{(L)} \, \vfield^{\mathrm{s},(L)}_T
  \label{eq:cross-layer}
\end{equation}
where $\sigma(w_g)$ is a learned sigmoid gate and $P^{(L)}$ is a
linear projection.

This is \textbf{iteration of the RG transformation}.  Each layer
receives the coarse-grained state from the layer below and operates
on it at a yet coarser scale.  With $N$ layers, the effective
temporal horizon grows as $\tau_{\mathrm{eff}} \sim \tau_s^N$---each
layer extends the memory exponentially.

\paragraph{The full system.}  Both fields contribute to key
modulation:
\begin{equation}
  \K_j'^{(h,L)} = \K_j^{(h,L)}
    + W_f^{(h,L)} \vfield^{\mathrm{f},(L)}_j
    + W_s^{(h,L)} \vfield^{\mathrm{s},(L)}_j
  \label{eq:full-key-mod}
\end{equation}

The attention mechanism then operates on the modulated keys:
\begin{equation}
  \ell_{ij}^{(h,L)}
    = \frac{\q_i^{(h,L)\top} \K_j'^{(h,L)}}{\sqrt{d_h}}
  \label{eq:full-attention}
\end{equation}

No scalar logit bias appears.  The field's influence is entirely
through key geometry.

\subsection{RG Interpretation: Why This Should Give Length Generalization}
\label{sec:rg-interpretation}

The connection between our architecture and the RG is not merely
analogical.  We identify three specific mechanisms:

\paragraph{Scale invariance from critical dynamics.}  The decay
rates $\retain_f^{i-j}$ and $\retain_s^{i-j}$ depend only on
the \emph{distance} $i - j$, not on absolute position.  The field's
memory horizon is position-independent.  At the optimal $\evap$
(the ``critical point''), the system's behavior is independent of
total sequence length---precisely the definition of length
generalization.

\paragraph{Relevant operators from the ablation.}  The
all-or-nothing ablation (\cref{sec:scalar-ablation}) maps directly
onto the RG concept of relevant operators:
\begin{itemize}
  \item \textbf{Feedback deposits} = relevant.  Without them, the
    field is a deterministic average (no state tracking).  Removing
    this operator drives the system to the trivial fixed point
    (chance accuracy).
  \item \textbf{Scale separation} = relevant.  Without it, there is
    no coarse-graining hierarchy.  A single decay rate means a single
    ``magnification''---no RG flow.
  \item \textbf{Cross-layer coupling} = relevant.  Without it, the
    RG transformation is applied once (within a layer) but never
    iterated.  The hierarchy has depth 1.
\end{itemize}
All five partial configurations are missing at least one relevant
operator and flow to the trivial fixed point.  Only the full system
sits near the nontrivial critical point.

\paragraph{Universality prediction.}  If the analogy holds, systems
with the same ``state complexity'' (number of bits tracked) should
exhibit the same generalization behavior regardless of task details.
We observe this: parity (1-bit) and cumsum mod 4 (2-bit) both show
monotonic degradation with length, with the field's benefit scaling
with available state capacity.

\paragraph{A prediction.}  The rescaling factor
$b = \tau_s/\tau_f$ should matter.  If $b$ is too small, the fast
and slow fields are not well-separated and the ``coarse-graining''
is trivial.  If $b$ is too large, information is lost in the
averaging.  There should be an optimal $b$---a critical ratio---that
maximizes generalization.  The sharp $\evap$ sensitivity we observe
($\evap = 0.05$ works, $\evap = 0.10$ collapses on cumsum) is
consistent with this.


% ══════════════════════════════════════════════════════════════
\section{Evidence I: Scalar Field Ablation}
\label{sec:scalar-ablation}

We first establish the all-or-nothing structure using scalar fields,
where the small parameter count makes controlled ablation
straightforward.

\subsection{Setup}

All transformer models: 4 layers, 4 heads, $d_{\mathrm{model}} = 64$,
ReLU MLP with expansion factor 4, RMS normalisation, learned
positional embeddings.  Total: $\sim$209K parameters across all
configurations.

Task: running parity.  Train on $n = 16$ examples ($T = 34$ tokens),
test at $n \in \{16, 32, 48, 64\}$.  3 seeds per configuration.

Six configurations, differing only in field mechanism:

\begin{table}[h]
\centering
\caption{Scalar field configurations.  All share the same transformer
  backbone; parameter counts are matched ($\pm$0.5\%).}
\label{tab:scalar-configs}
\begin{tabular}{@{}lccccr@{}}
\toprule
Config & Field & Feedback & Multi-scale & Cross-layer & Params \\
\midrule
\texttt{no\_field}     & --- & --- & --- & --- & 208,512 \\
\texttt{input\_field}  & $\checkmark$ & & & & 209,536 \\
\texttt{feedback\_field} & $\checkmark$ & $\checkmark$ & & & 209,536 \\
\texttt{multiscale\_field} & $\checkmark$ & & $\checkmark$ & & 209,616 \\
\texttt{hierarchical\_field} & $\checkmark$ & & $\checkmark$ & $\checkmark$ & 209,664 \\
\texttt{hierarchical\_feedback} & $\checkmark$ & $\checkmark$ & $\checkmark$ & $\checkmark$ & 209,664 \\
\bottomrule
\end{tabular}
\end{table}

\subsection{Results: All-or-Nothing}

\begin{table}[h]
\centering
\caption{Scalar field ablation on parity.  Chance is 0.50.
  Only the full system generalises.  Mean $\pm$ std over 3 seeds.}
\label{tab:scalar-results}
\begin{tabular}{@{}l cccc@{}}
\toprule
& \multicolumn{4}{c}{Accuracy at $n$ examples ($T$ tokens)} \\
\cmidrule(l){2-5}
Config & $n{=}16$ & $n{=}32$ & $n{=}48$ & $n{=}64$ \\
\midrule
\texttt{no\_field}        & $1.000 \pm .000$ & $0.589 \pm .037$ & $0.532 \pm .021$ & $0.523 \pm .019$ \\
\texttt{input\_field}     & $1.000 \pm .000$ & $0.488 \pm .008$ & $0.498 \pm .002$ & $0.496 \pm .005$ \\
\texttt{feedback\_field}  & $1.000 \pm .000$ & $0.517 \pm .012$ & $0.503 \pm .004$ & $0.508 \pm .010$ \\
\texttt{multiscale}       & $1.000 \pm .000$ & $0.496 \pm .008$ & $0.499 \pm .002$ & $0.495 \pm .002$ \\
\texttt{hierarchical}     & $1.000 \pm .000$ & $0.483 \pm .000$ & $0.503 \pm .006$ & $0.498 \pm .006$ \\
\texttt{hier.\ feedback}  & $\mathbf{1.000 \pm .000}$ & $\mathbf{0.808 \pm .031}$ & $\mathbf{0.710 \pm .110}$ & $\mathbf{0.653 \pm .079}$ \\
\bottomrule
\end{tabular}
\end{table}

The result is binary.  Five configurations are indistinguishable from
chance at all OOD lengths.  The sixth---and \emph{only} the
sixth---generalises.  In RG language: five configurations are missing
a relevant operator and flow to the trivial fixed point.  The full
system sits near the nontrivial critical point.

\paragraph{Why each component fails alone.}
\begin{itemize}
  \item \textbf{Input deposits}: the field is a deterministic
    function of input---a moving average that provides no state.
  \item \textbf{Feedback without multi-scale}: single timescale
    cannot simultaneously provide local precision and long-range
    memory.  No scale separation $\Rightarrow$ no coarse-graining.
  \item \textbf{Multi-scale without feedback}: two moving averages
    are still averages.  The coarse-graining has nothing useful to
    coarse-grain.
  \item \textbf{Cross-layer without feedback}: vertical plumbing
    transmits noise.  The RG iteration operates on empty state.
\end{itemize}

\subsection{SSM Baselines}

\begin{table}[h]
\centering
\caption{Dedicated recurrent models achieve perfect generalization,
  confirming that recurrence trivially solves parity.}
\label{tab:ssm-results}
\begin{tabular}{@{}l rcccc@{}}
\toprule
Model & Params & $n{=}16$ & $n{=}32$ & $n{=}48$ & $n{=}64$ \\
\midrule
S4D              & 167K & $1.000$ & $1.000$ & $1.000$ & $1.000$ \\
Selective SSM    & 184K & $1.000$ & $1.000$ & $1.000$ & $1.000$ \\
GRU              & 227K & $1.000$ & $1.000$ & $1.000$ & $1.000$ \\
\midrule
\textit{Hier.\ feedback (scalar)} & \textit{210K} & \textit{1.000} & \textit{0.808} & \textit{0.710} & \textit{0.653} \\
\bottomrule
\end{tabular}
\end{table}

SSMs achieve perfect generalization because they \emph{are} the fixed
point: pure recurrence is inherently scale-invariant for state
tracking.  The question is whether we can approach this fixed point
from the attention side.


% ══════════════════════════════════════════════════════════════
\section{Evidence II: Vector Fields and State Capacity}
\label{sec:vector-results}

The scalar field's limitation is capacity: $H$ scalars provide
$O(H)$ bits of state.  Upgrading to vector fields tests whether the
RG hierarchy works with richer ``spins.''

Results are from a companion dual-channel experiment using 2 layers,
4 heads, $d_{\mathrm{model}} = 48$, $d_f = 12$ ($\sim$60K params),
RoPE positional encoding, trained on $n = 16$ ($T = 34$).

\begin{table}[h]
\centering
\caption{Vector social field (additive key modulation) vs.\ scalar
  field and baseline.  The vector field's advantage grows
  dramatically with state complexity.}
\label{tab:vector-results}
\begin{tabular}{@{}lcccccc@{}}
\toprule
& \multicolumn{3}{c}{\textbf{Parity (1-bit)}}
& \multicolumn{3}{c}{\textbf{Cumsum mod 16 (4-bit)}} \\
\cmidrule(lr){2-4} \cmidrule(lr){5-7}
Mode & T=34 & T=66 & T=98
     & T=34 & T=66 & T=98 \\
\midrule
Baseline
  & $1.000$ & $0.995$ & $0.525$
  & $1.000$ & $0.279$ & $0.062$ \\
Scalar feedback
  & $1.000$ & $0.949$ & $0.699$
  & $1.000$ & $0.870$ & $0.216$ \\
\textbf{Vector additive}
  & $1.000$ & $0.979$ & $\mathbf{0.738}$
  & $1.000$ & $\mathbf{0.999}$ & $\mathbf{0.419}$ \\
\bottomrule
\end{tabular}
\end{table}

\paragraph{Key observations.}
\begin{itemize}
  \item On cumsum mod 16 (4-bit state) at $2\times$ training length:
    vector additive achieves \textbf{0.999} vs.\ baseline
    \textbf{0.279}---a $+0.721$ improvement.  This is the largest
    effect in any experiment.  The vector field's $H \times d_f = 48$
    values of recurrent state can encode the 4-bit running sum that
    scalars cannot.
  \item The advantage grows with state complexity: $+0.039$ on parity
    (1-bit) vs.\ $+0.203$ on cumsum (4-bit) at $3\times$.  This is
    the capacity scaling predicted by the RG interpretation: richer
    ``spins'' (vectors vs.\ scalars) carry more information through
    the coarse-graining hierarchy.
  \item Additive key modulation outperforms multiplicative gating
    (not shown; see companion paper).  The additive mechanism has a
    natural geometric interpretation: the field \emph{translates}
    keys in a learned direction, preserving the attention mechanism's
    dot-product structure.
\end{itemize}

\subsection{Multi-Task State-Capacity Scaling (Scalar Fields)}

\begin{table}[h]
\centering
\caption{State-capacity scaling at $4\times$ training length.
  The field's benefit degrades monotonically with state complexity.}
\label{tab:multitask}
\begin{tabular}{@{}l cccc@{}}
\toprule
& Parity & Cumsum mod 4 & Cumsum mod 16 & Flipflop \\
& (1-bit) & (2-bit) & (4-bit) & (gating) \\
\midrule
No field       & 0.52 & 0.25 & 0.06 & 0.10 \\
Hier.\ feedback & \textbf{0.65} & \textbf{0.40} & \textbf{0.15} & 0.15 \\
S4D            & 1.00 & 1.00 & 0.95 & 0.20 \\
\midrule
Random         & 0.50 & 0.25 & 0.06 & 0.06 \\
\bottomrule
\end{tabular}
\end{table}

The monotonic degradation from parity to cumsum mod 16 confirms the
$O(H)$-bit capacity prediction for scalar fields.  The flipflop
result reveals a boundary: the field helps with \emph{accumulation}
tasks (where all inputs contribute) but not \emph{gating} tasks
(selective state update).  In RG terms, gating requires a different
universality class---one where the coarse-graining rule is
input-dependent, not fixed.


% ══════════════════════════════════════════════════════════════
\section{The Renormalization Picture}
\label{sec:rg-analysis}

\subsection{Block-Spin Analogy Made Precise}

In Kadanoff's block-spin construction, one partitions a
$d$-dimensional lattice into blocks of linear size $b$, replaces each
block's spins with a single ``block spin'' (typically the block
average or majority), and rescales the lattice so the block lattice
has the same spacing as the original.

Our temporal analogue:

\begin{center}
\begin{tabular}{@{}lll@{}}
\toprule
\textbf{Block-spin RG} & \textbf{Our hierarchy} & \textbf{Component} \\
\midrule
Lattice sites & Sequence positions & Input tokens \\
Spin at site $i$ & Key vector $\K_i$ & Level 0 \\
Block of size $b$ & Window of $\tau_f$ tokens & Fast-field kernel \\
Block spin & Fast-field vector $\vfield^{\mathrm{f}}_i$ & Level 1 \\
Super-block & Window of $\tau_s$ tokens & Slow-field kernel \\
Super-block spin & Slow-field vector $\vfield^{\mathrm{s}}_i$ & Level 2 \\
RG iteration & Cross-layer coupling & Level 3 \\
\midrule
Rescaling factor $b$ & $\tau_s / \tau_f \approx 8$ & \\
Critical point & Optimal $\evap$ & Edge of chaos \\
Fixed-point condition & Length generalization & Scale invariance \\
\bottomrule
\end{tabular}
\end{center}

The exponential kernel $\retain^{i-j}$ is not the only possible
coarse-graining rule, but it has a key property: \textbf{it is
translation-invariant}.  The coarse-graining at position $i$ uses the
same kernel as at position $i + \Delta$, regardless of $\Delta$.
This is the structural reason the field generalizes across lengths.

\subsection{Scale Separation and the Adiabatic Condition}

For the hierarchy to function as a proper RG, the scales must be
well-separated: $\tau_s \gg \tau_f$.  In our experiments,
$\tau_s/\tau_f \approx 8$, which provides roughly one order of
magnitude of separation.

The separation enables an \textbf{adiabatic approximation}: from the
perspective of the fast field, the slow field is approximately
constant; from the perspective of the slow field, the fast field
fluctuations are averaged out.  This means:
\begin{itemize}
  \item The fast field tracks local structure without interference
    from long-range memory.
  \item The slow field accumulates global structure without being
    disrupted by token-level noise.
  \item Each scale maintains its own dynamics while exchanging
    information through the hierarchy.
\end{itemize}

The ant-trail analogy: trail pheromone (fast, volatile) guides local
foraging; colony pheromone (slow, persistent) encodes the nest
location.  They operate at different timescales and influence
different behaviors, coupled only through the ants' movements.

\subsection{Why the Single-Scale System Fails: No RG Flow}

A single decay rate $\evap$ gives a single timescale $\tau = 1/\evap$.
This is a system with one ``magnification''---there is no RG flow
because there is no hierarchy to flow through.  The system either
remembers (if $\evap$ is small enough for the test length) or
forgets (if $\evap$ is too large), with no mechanism to bridge
scales.

The failure of single-scale feedback (\cref{tab:scalar-results}) is
not a failure of recurrence but a failure of \emph{scale structure}.
The system has state but no mechanism to organize it across temporal
scales.

\subsection{Evaporation Rate as a Control Parameter}

In the RG framework, a \emph{control parameter} tunes the system
through a phase transition.  The evaporation rate $\evap$ plays this
role:

\begin{itemize}
  \item $\evap \to 1$ (rapid erosion): the field decays instantly.
    No memory.  The system is in the ``disordered phase''---attention
    operates as if no field exists.
  \item $\evap \to 0$ (no erosion): the field saturates.  All
    positions accumulate equal weight.  The system is in the
    ``ordered phase''---attention is dominated by the field, losing
    content sensitivity.
  \item $\evap$ at criticality: the field maintains structured,
    content-dependent memory that neither saturates nor vanishes.
    The system sits at the edge of chaos, maximizing information
    transmission across scales.
\end{itemize}

The sharp $\evap$ sensitivity observed empirically ($\evap = 0.05$
works on cumsum; $\evap = 0.10$ collapses to 0.256) is the signature
of a narrow critical window---a hallmark of systems near a phase
transition.


% ══════════════════════════════════════════════════════════════
\section{The Multiscale Social Field}
\label{sec:combined}

The preceding sections present two evidence streams from separate
toy-problem experiments: Section~\ref{sec:scalar-ablation} shows that
multi-scale structure is necessary (using scalar fields), and
Section~\ref{sec:vector-results} shows that vector social fields
provide the state capacity (using a single timescale).  We now
describe the architecture that combines both: the \emph{multiscale
social field}---a vector social field with two decay timescales
organized as a renormalization hierarchy.

\begin{enumerate}
  \item \textbf{Vector social fields}
    (\cref{sec:vector-results}) provide the state capacity.  Each
    head carries $d_f$-dimensional recurrent state through additive
    key modulation.
  \item \textbf{Multi-scale structure}
    (\cref{sec:scalar-ablation}) provides the temporal hierarchy.
    Fast and slow fields implement a two-level block-spin
    coarse-graining.
  \item \textbf{Cross-layer coupling} iterates the RG
    transformation, extending the effective memory exponentially with
    depth.
  \item \textbf{Feedback deposits} make the field content-dependent,
    creating genuine state rather than deterministic averages.
\end{enumerate}

The combined system for layer $L$, head $h$, position $i$:
\begin{align}
  \vfield^{\mathrm{f},(L,h)}_i &= \retain_f \cdot
    \vfield^{\mathrm{f},(L,h)}_{i-1}
    + W_{\mathrm{dep}}^{(L,h)} \, \mathbf{o}_i^{(L,h)}
    \label{eq:combined-fast} \\
  \vfield^{\mathrm{s},(L,h)}_i &= \retain_s \cdot
    \vfield^{\mathrm{s},(L,h)}_{i-1}
    + h^{(L,h)}\!\left(\vfield^{\mathrm{f},(L,h)}_i\right)
    \label{eq:combined-slow} \\
  \vfield^{\mathrm{s},(L,h)}_0 &= \sigma\!\left(w_g^{(L-1,h)}\right)
    \cdot P^{(L-1,h)} \, \vfield^{\mathrm{s},(L-1,h)}_T
    \label{eq:combined-cross} \\
  \K_j'^{(L,h)} &= \K_j^{(L,h)}
    + W_f^{(L,h)} \vfield^{\mathrm{f},(L,h)}_j
    + W_s^{(L,h)} \vfield^{\mathrm{s},(L,h)}_j
    \label{eq:combined-keys}
\end{align}

\paragraph{Scaling.}  The field adds $O(H \cdot d_f)$ parameters per
layer for the deposit, read, and projection matrices.  With
$d_f = d_h$, this is $<1\%$ parameter overhead at typical model sizes.
The recurrent state per position is $2 \cdot H \cdot d_f$ values
(fast + slow per head)---far richer than scalar fields but far
smaller than a full SSM hidden state.

\paragraph{Predictions.}  Combining the scalar ablation
(\cref{tab:scalar-results}) with the vector capacity results
(\cref{tab:vector-results}), this architecture predicts:
\begin{itemize}
  \item The hierarchy must be complete---partial configurations
    should fail.  \emph{Confirmed} on natural language
    (\cref{sec:natural-language}): single-field variants show
    2.18\% ablation delta vs.\ 36.6\% for multiscale.
  \item The model should spontaneously discover scale separation
    when given learnable retention.  \emph{Confirmed}: the slow
    field converges to $\retain_s \approx 0.994$ (half-life
    $\sim$120 tokens) while the fast field retains
    $\retain_f = 0.95$ (half-life $\sim$13 tokens).
  \item The field's contribution should grow monotonically with
    position.  \emph{Confirmed}: the per-chunk gradient is
    $0\% \to 0.7\% \to 8.3\% \to 36.6\%$.
  \item The multiscale architecture should improve perplexity on
    natural language at scale.  \emph{Confirmed}: 19\% PPL reduction
    on OWT at 137.8M parameters (\cref{sec:owt-multiscale}).
  \item The architecture should close the gap to SSMs on
    multi-bit state-tracking tasks.  \emph{Untested} at scale.
  \item Sharp $\evap$ sensitivity should persist.
    \emph{Untested}---the TinyStories sweep tested architectural
    variants, not $\evap$ sweeps.
  \item Multiple layers should develop different timescales.
    \emph{Partially confirmed, mechanism revised}: on TinyStories
    (4~layers), all layers converge to similar long half-lives
    ($\sim$200 tokens), and the slow field is active everywhere.  On
    OWT (12~layers), only Layer~0 retains a long half-life (300
    tokens); deeper layers deactivate the slow field
    (\cref{sec:owt-multiscale}).  A retention strategy sweep
    (\cref{sec:retention-sweep}) shows that \emph{forcing} a
    timescale hierarchy hurts performance---the model adaptively
    allocates slow-field retention where most needed rather than
    building a fixed cross-layer hierarchy.
\end{itemize}


% ══════════════════════════════════════════════════════════════
\section{Empirical Validation: Multiscale Fields on Natural Language}
\label{sec:natural-language}

The preceding experiments test the field on synthetic state-tracking
tasks designed to isolate its capabilities.  A critical question
remains: does the multi-scale structure matter when the field is
embedded in a language model trained on natural text?

\subsection{Setup}

We train five architectural variants on TinyStories (Eldan \& Li,
2023)---2.2M synthetic children's stories---using identical
hyperparameters: $d_{\mathrm{model}} = 192$, 4 heads, 4 layers,
context length 512, chunk size 128, 3000 training steps.  All
variants use the chunk-parallel vector field with additive key
modulation, differing only in field configuration:

\begin{table}[h]
\centering
\caption{Architectural variants tested on TinyStories.}
\label{tab:variants}
\begin{tabular}{@{}lrcc@{}}
\toprule
Variant & Params & Field configuration & Retention \\
\midrule
\texttt{baseline}   & 11.42M & No field & --- \\
\texttt{fixed}      & 11.58M & Single field, $\evap = 0.05$ & Fixed, $\retain = 0.95$ \\
\texttt{learnable}  & 11.58M & Single field, learnable $\evap$ & Per-head sigmoid \\
\texttt{multiscale} & 11.73M & Fast ($\evap = 0.05$) + slow (learnable) & Fixed + per-head \\
\texttt{crosslayer} & 11.58M & Single field, vertical coupling & Fixed, $\retain = 0.95$ \\
\bottomrule
\end{tabular}
\end{table}

\subsection{Evaluation: Per-Chunk Ablation}

Standard perplexity averages over all positions equally, obscuring
position-dependent effects.  The field modulates keys starting from
chunk 1 (positions 128+), so its contribution should be zero in
chunk 0 and grow in later chunks.  We measure this by zeroing
$W_{\mathrm{mod}}$ (and $W_{\mathrm{mod,slow}}$ for multiscale) and
computing per-chunk perplexity change.

\subsection{Results}

\begin{table}[h]
\centering
\caption{Per-chunk ablation: perplexity increase (\%) when field
  modulation is zeroed out.  The signature of useful field state
  is growing $\Delta$\% at later positions.}
\label{tab:ablation}
\begin{tabular}{@{}lccccr@{}}
\toprule
& Val PPL & Chunk 0 & Chunk 1 & Chunk 2 & Chunk 3 \\
& & (0--127) & (128--255) & (256--383) & (384--511) \\
\midrule
\texttt{baseline}   & 15.6 & --- & --- & --- & --- \\
\texttt{fixed}      & 15.6 & $+0.00$\% & $+0.05$\% & $+0.35$\% & $+2.18$\% \\
\texttt{learnable}  & 15.6 & $+0.00$\% & $+0.05$\% & $+0.35$\% & $+2.18$\% \\
\texttt{multiscale} & \textbf{15.4} & $+0.00$\% & $+0.70$\% & $+8.28$\% & $\mathbf{+36.60}$\% \\
\texttt{crosslayer} & 15.6 & $+0.00$\% & $+0.05$\% & $+0.35$\% & $+2.18$\% \\
\bottomrule
\end{tabular}
\end{table}

\paragraph{Multiscale dominates.}  The multiscale variant is the only
configuration that achieves both better aggregate perplexity (15.4 vs
15.6) and a dramatically larger ablation delta.  Removing the field
causes a \textbf{36.6\% perplexity increase} at positions 384--511,
compared to 2.18\% for all single-field variants.  The positional
gradient---$0\% \to 0.7\% \to 8.3\% \to 36.6\%$---is exactly the
signature predicted by the RG interpretation: accumulated cross-chunk
state becomes increasingly essential at later positions.

\paragraph{Learnable retention converges back to the init.}
All heads converged to $\retain \approx 0.948$ (init was 0.95),
with half-lives of $\sim$13 tokens across all layers.  With a single
field, the model found no reason to lengthen the timescale---the
fast decay is locally optimal for immediate chunk-to-chunk modulation.
This confirms that the limitation is \emph{architectural} (single
timescale), not a matter of finding the right decay rate.

\paragraph{The slow field learned long-range retention.}
The multiscale variant's slow field converged to
$\retain_s \approx 0.994$ (half-life $\sim$115--127 tokens),
spanning most of the 512-token context window.  The fast field
retained its default decay.  The model spontaneously discovered
scale separation: a fast field for local dynamics and a slow field
for long-range memory---precisely the hierarchy the RG theory
prescribes.

\paragraph{Cross-layer coupling had no effect.}  Passing field state
vertically between layers produced identical results to the fixed
single-field variant.  The bottleneck was \emph{temporal range} (how
far back the field remembers), not \emph{depth} (how many layers
process the field).  Cross-layer coupling cannot compensate for a
single timescale that forgets within one chunk.

\subsection{Concurrent Validation: OpenWebText 125M}

In parallel, we trained a 130.7M parameter GPT (GPT-2 Small scale)
with a single fixed field ($\evap = 0.05$) on OpenWebText (9B tokens,
context 1024, chunk 128).  Per-chunk ablation at two checkpoints
confirmed the field's contribution grows with training:

\begin{table}[h]
\centering
\caption{OWT per-chunk ablation delta at positions 896--1023
  (last chunk) over training.}
\label{tab:owt-trend}
\begin{tabular}{@{}lccc@{}}
\toprule
Checkpoint & Last-chunk $\Delta$\% & Val PPL & Training \% \\
\midrule
Step 30K (single-field) & $+0.45$\% & 30.4 & 22\% \\
Step 60K (single-field) & $+0.60$\% & 27.4 & 44\% \\
\midrule
Step 70K (multiscale) & $+0.86$\% & \textbf{22.2} & 100\% \\
\bottomrule
\end{tabular}
\end{table}

The single-field OWT model shows a growing but small late-chunk
signal, consistent with our TinyStories finding: the fixed
$\evap = 0.05$ field (half-life $\sim$13 tokens) contributes
measurably at long range but is architecturally limited.  With
$\retain^{128} \approx 0.001$, deposits from one chunk are
essentially gone by the next.

\subsection{Full-Scale Validation: Multiscale on OpenWebText}
\label{sec:owt-multiscale}

We warm-started a multiscale variant from the step-60K single-field
checkpoint, loading all 111 matched parameters and initializing 36
new slow-field parameters fresh (3 per layer: \texttt{slow\_retain\_logit},
$W_{\mathrm{dep,slow}}$, $W_{\mathrm{mod,slow}}$).  The multiscale
model (137.8M parameters, $<0.1\%$ overhead) was trained for 70K
additional steps on OpenWebText with the same hyperparameters
($d = 768$, 12 heads, 12 layers, context 1024, chunk 128).

\begin{table}[h]
\centering
\caption{OWT validation perplexity: single-field vs.\ multiscale.
  The multiscale variant surpasses the single-field baseline and
  continues improving throughout training.}
\label{tab:owt-multiscale-ppl}
\begin{tabular}{@{}lccc@{}}
\toprule
Model & Steps trained & Val PPL & Val loss \\
\midrule
Single-field (fixed $\evap = 0.05$) & 60K & 27.4 & 3.31 \\
\textbf{Multiscale} (warm-started) & +70K & \textbf{22.2} & \textbf{3.10} \\
\bottomrule
\end{tabular}
\end{table}

The multiscale variant achieves a \textbf{19\% perplexity reduction}
(27.4 $\to$ 22.2), surpassing the single-field model within the first
28K steps and continuing to improve throughout training.

\paragraph{Per-chunk ablation confirms field utility.}
At step 25K, we probed the slow field's contribution by zeroing
$W_{\mathrm{mod,slow}}$ and measuring per-chunk loss change:

\begin{table}[h]
\centering
\caption{Per-chunk ablation at step 25K of multiscale OWT training.
  Positive delta means the field helps prediction at that position.
  ``Slow share'' is the fraction of total field effect attributable
  to the slow field alone.}
\label{tab:owt-chunk-ablation}
\begin{tabular}{@{}lccccc@{}}
\toprule
& Chunk 0 & Chunk 3 & Chunk 5 & Chunk 6 & Chunk 7 \\
& (0--127) & (384--511) & (640--767) & (768--895) & (896--1023) \\
\midrule
Slow $\Delta$ & $+0.000$ & $+0.001$ & $+0.003$ & $+0.004$ & $+0.003$ \\
All fields $\Delta$ & $+0.000$ & $+0.001$ & $+0.006$ & $+0.006$ & $+0.009$ \\
Slow share & --- & 85.8\% & 56.0\% & 63.0\% & 36.7\% \\
\bottomrule
\end{tabular}
\end{table}

The slow field accounts for \textbf{37--97\% of the total field effect}
in the second half of the sequence (chunks 3--7), confirming that the
learnable slow component is not dead weight.  The overall pattern
replicates the TinyStories finding: ablation delta grows monotonically
with position, and the slow field is responsible for the majority of
the effect in middle chunks.

\paragraph{Emergent layer specialization.}
The most striking result is \emph{how} the model organized its
slow-field retention across layers.  All 12 layers were initialized
at half-life $\approx 200$ tokens.  By the end of training:

\begin{table}[h]
\centering
\caption{Learned slow-field retention at end of training (step 70K).
  The model self-organized into one long-memory layer (Layer~0) and
  11 short-memory layers.}
\label{tab:owt-retention}
\begin{tabular}{@{}lcc@{}}
\toprule
Layer & Retention ($\retain_s$) & Half-life (tokens) \\
\midrule
Layer 0 & 0.9967 & 300 \\
Layer 1 & 0.6580 & 2 \\
Layer 2 & 0.6242 & 2 \\
Layers 3--8 & 0.15--0.47 & $<$1 \\
Layers 9--11 & 0.11--0.12 & $<$1 \\
\bottomrule
\end{tabular}
\end{table}

The model learned a \textbf{single long-memory gateway}: Layer~0's
slow field acts as a persistent 300-token memory, while all deeper
layers collapsed their slow fields to sub-token half-lives---effectively
disabling them.  This is a parsimonious solution: Layer~0 captures
long-range context at the bottom of the residual stream, and deeper
layers access this information through the residual connection rather
than maintaining their own slow fields.

\paragraph{Interpretation.}  This result both confirms and
complicates the RG theory.  The model \emph{did} discover that
long-range memory improves prediction (the PPL drop and ablation
deltas are unambiguous).  But it did \emph{not} build a distributed
hierarchy of timescales across layers---it found that a single
long-memory entry point is sufficient, at least at this scale.

There are two possible readings: (1) a single gateway genuinely
suffices for 1024-token contexts, and the distributed hierarchy
would only emerge at longer contexts where residual propagation
from Layer~0 becomes lossy; or (2) the optimizer dynamics cause
premature collapse---the gradient signal for long-range dependencies
is weak relative to local prediction, and high learning rates kill
the slow field in deeper layers before it can prove useful.

We note that the half-life of Layer~0 grew monotonically throughout
training (200 $\to$ 261 $\to$ 300 tokens), suggesting the model
would continue pushing for longer retention given more training or
longer contexts.  To disambiguate the two readings---genuine
preference vs.\ optimizer artifact---we ran a controlled sweep of
retention training strategies.

\subsection{Retention Strategy Sweep}
\label{sec:retention-sweep}

We tested four retention strategies on TinyStories ($d = 128$,
4 heads, 4 layers, context 512, chunk 128, 5000 steps) to determine
whether the OWT layer collapse was an optimizer artifact or a genuine
architectural preference:

\begin{table}[h]
\centering
\caption{Retention strategy sweep.  ``Last-chunk $\Delta$'' is the
  loss increase when zeroing slow-field modulation at positions
  384--511.  Larger $\Delta$ means the slow field is more important.}
\label{tab:retention-sweep}
\begin{tabular}{@{}lccl@{}}
\toprule
Variant & Val PPL & Last-chunk $\Delta$ & Learned half-lives \\
\midrule
\texttt{ms\_uniform}
  & \textbf{18.1} & $\mathbf{+0.397}$
  & L0:220, L1:213, L2:211, L3:212 \\
\texttt{ms\_low\_lr}
  & 18.5 & $+0.391$
  & L0:202, L1:201, L2:201, L3:201 \\
\texttt{ms\_fixed\_hier}
  & 18.2 & $+0.178$
  & L0:500*, L1:171*, L2:58*, L3:20* \\
\texttt{ms\_diverse\_init}
  & 18.2 & $+0.187$
  & L0:528, L1:186, L2:68, L3:22 \\
\bottomrule
\multicolumn{4}{@{}l}{\small *Frozen (not learned).}
\end{tabular}
\end{table}

\paragraph{Uniform retention wins decisively.}
The default multiscale variant with uniform initialization and
standard learning rate produced the best PPL (18.1) and the
largest ablation delta (+0.397)---more than double the hierarchical
variants (+0.178, +0.187).  The per-chunk progression is
$+0.000 \to +0.002 \to +0.081 \to +0.397$, confirming the slow
field's contribution grows monotonically with position.

\paragraph{Forced hierarchy hurts.}  Both hierarchical
variants---one with frozen geometric spacing (500$\to$171$\to$58$\to$20),
one with the same initialization but learnable---performed
identically (PPL 18.2) and had \emph{half} the ablation signal.  The
problem is structural: Layer~3 at half-life 20 tokens ($\retain^{128}
\approx 0.002$) loses its deposits within a single 128-token chunk.
The near-uniform half-life of $\sim$200--220 tokens is better matched
to the chunk structure.

\paragraph{Diverse initialization is sticky.}  The learnable
variant initialized at the same geometric spread converged to
L0:528, L1:186, L2:68, L3:22---barely moved from initialization.
The optimizer did not collapse the hierarchy toward uniformity; it
preserved the initial spread.  Yet this preserved hierarchy
performed worse than the uniform variant.  This rules out the
hypothesis that the hierarchy is better but hard to discover: when
given a hierarchy, the model kept it and still underperformed.

\paragraph{Lower learning rate is slightly worse.}  Reducing the
retain logit LR by $10\times$ prevented adaptation (all layers stayed
near the hl$\approx$200 initialization) and slightly hurt PPL (18.5).

\paragraph{Interpretation.}  These results resolve the ambiguity
from the OWT experiment---but in a subtler way than simple
``collapse.''  On TinyStories (4 layers), the winning configuration
uses the slow field \emph{everywhere}: all layers converge to
half-lives of 211--220 tokens and the ablation signal is large.  On
OWT (12 layers), the model concentrates slow-field retention in
Layer~0 and lets deeper layers rely on the richer attention pathway
available at 137.8M parameters.

The multiscale field is therefore \textbf{adaptive}: the architecture
provides the capability at every layer, and the optimizer activates
it where most efficient.  In a small model with limited attention
capacity, every layer benefits from persistent state.  In a larger
model, the residual stream carries Layer~0's field-enriched
representations upward, and deeper layers can free their slow fields
without losing the information.

The key insight is that the \textbf{modulation weights}
($W_{\mathrm{dep,slow}}$, $W_{\mathrm{mod,slow}}$), not the
retention values, do the heavy lifting.  The retention just needs to
be long enough to carry state across chunks; beyond that, the
deposit and read projections determine what is stored and retrieved.
When forced into a fixed hierarchy, the short-retention layers
(Layer~3 at hl=20) lose deposits within a single chunk, wasting
capacity on a timescale the chunk structure cannot exploit.

This refines the RG interpretation: the fast/slow split within each
layer provides the essential scale separation.  The cross-layer
dimension is not a fixed hierarchy but an \emph{adaptive allocation}
of a shared resource---long-range memory---to whatever depth the
model and data require.


% ══════════════════════════════════════════════════════════════
\section{Discussion}
\label{sec:discussion}

\subsection{Where This Sits in the Design Space}

\begin{center}
\begin{tabular}{@{}lccc@{}}
\toprule
& \textbf{Attention} & \textbf{This work} & \textbf{SSM} \\
\midrule
Random access & Full & Full & None \\
Recurrent state & None & $2Hd_f$ per pos. & Full \\
Length generalization & None & Partial--good & Perfect \\
Parallelizable & Fully & Partially & Via scan \\
\bottomrule
\end{tabular}
\end{center}

The architecture does not replace attention with recurrence (SSMs)
or alternate between them (Jamba/Griffin).  It augments attention
with a lightweight recurrence---a vector trail---organized as a
renormalization hierarchy.  The attention mechanism remains primary;
the field is a ``guide'' that shapes what attention looks for based
on accumulated history.

\subsection{Not Trying to Beat SSMs}

SSMs are the natural solution when full recurrence is needed.  Our
goal is different: to give transformers---which dominate production
deployments---a \emph{biologically motivated} form of sequence
memory.  The question is not ``can this beat Mamba on parity?'' (it
cannot) but ``can the RG hierarchy give attention a principled,
interpretable form of state tracking that standard attention lacks?''

The vector field results on cumsum mod 16 (0.999 at $2\times$ vs.\
0.279 baseline) suggest the answer is yes for tasks within the
field's capacity.  The renormalization framing provides hyperparameter
guidance ($\evap$ controls a phase transition) and architectural
guidance (the hierarchy must be complete---no shortcuts).

\subsection{Biological Plausibility and Interpretability}

The architecture has a direct biological interpretation:
\begin{itemize}
  \item Fast field = trail pheromone (volatile, local)
  \item Slow field = colony pheromone (persistent, global)
  \item Cross-layer = information transfer between organizational
    levels (forager $\to$ colony)
  \item Key modulation = pheromone reshaping the ``search landscape''
\end{itemize}

This is interpretable in a way that SSM hidden states are not.  One
can inspect what direction the field is shifting keys, at what
timescale, and how it flows between layers.  The field is a
``social memory'' that the model deposits and reads---a form of
external working memory embedded in the attention mechanism.

\subsection{Limitations and Open Questions}

\begin{enumerate}
  \item \textbf{Adaptive rather than fixed hierarchy.}
    The RG theory predicts that different layers should adopt
    different timescales.  In practice, the model adaptively
    allocates slow-field retention: uniformly across all layers in
    small models (TinyStories), but concentrated in Layer~0 in larger
    models (OWT).  Forcing a fixed geometric hierarchy hurts
    performance (\cref{sec:retention-sweep}).  The cross-layer
    dimension is better understood as adaptive resource allocation
    than as a fixed RG iteration.  Whether deeper models at longer
    context lengths would re-activate slow fields in intermediate
    layers is unknown.
  \item \textbf{Scale beyond GPT-2 Small.}  Results span from
    $\sim$60K-parameter toy models to 137.8M parameters on
    OpenWebText.  Whether the multiscale advantage persists at
    GPT-2 Medium and larger is unknown.
  \item \textbf{Sequential cost.}  Feedback deposits require
    sequential processing ($O(T)$ bottleneck).  The chunk-parallel
    implementation amortizes this but still processes chunks serially.
  \item \textbf{Is the RG analogy predictive?}  The TinyStories
    results confirm two qualitative predictions: the hierarchy must
    be complete (multiscale $\gg$ single-field) and the model
    spontaneously discovers scale separation.  Quantitative
    predictions---scaling laws for $b = \tau_s/\tau_f$, critical
    exponents, universality---remain untested.
  \item \textbf{Continuous spectrum.}  Real RG systems have a
    continuous spectrum of scales.  Our two-level hierarchy
    (fast/slow) is the minimal realization.  Richer hierarchies
    (three or more timescales) may be needed for harder tasks
    or longer contexts.
\end{enumerate}


% ══════════════════════════════════════════════════════════════
\section{Conclusion}
\label{sec:conclusion}

Length generalization is a scale-invariance problem.  We have
argued---through controlled ablations on synthetic tasks, capacity
experiments, and preliminary natural language training---that a
renormalization hierarchy within transformer attention provides a
principled route to this invariance.

The hierarchy has three levels: a fast vector field capturing local
dynamics, a slow vector field performing temporal coarse-graining,
and cross-layer coupling iterating the transformation.  The synthetic
ablation reveals that these levels are \emph{relevant operators} in
the RG sense: removing any one drives the system to the trivial
fixed point on state-tracking tasks.

The scalar version of this hierarchy is clean but weak---it
generalizes partially on 1-bit tasks but cannot compete with
dedicated recurrent models.  The vector version, which modulates
attention geometry rather than logit magnitudes, provides
dramatically higher state capacity and closes much of the gap on
multi-bit tasks.

The TinyStories experiments (\cref{sec:natural-language}) provide the
first evidence that the hierarchy matters for natural language: the
multiscale variant produces a 36.6\% perplexity increase in the last
chunk when disabled, compared to 2.18\% for single-field variants.
The full-scale OWT experiment (\cref{sec:owt-multiscale}) confirms
this at the 137.8M-parameter scale: the multiscale model achieves a
\textbf{19\% perplexity reduction} (27.4 $\to$ 22.2) over the
single-field baseline, with per-chunk ablation showing the slow field
responsible for 37--97\% of the total field effect in later chunks.

The renormalization framing makes specific predictions confirmed by
these results: the hierarchy must be complete (multiscale $\gg$
single-field), the model should discover scale separation when given
the capacity (it does), and the field's contribution should grow with
position (the chunk-by-chunk gradient is monotonically increasing).

However, the predicted cross-layer timescale hierarchy does not
take the fixed form the RG analogy suggests.  On TinyStories
(4~layers), all layers activate the slow field at similar retention
($\sim$200 tokens).  On OWT (12~layers), the model concentrates
long-range memory in Layer~0 (half-life 300 tokens), deactivating
the slow field elsewhere.  A controlled retention strategy sweep
(\cref{sec:retention-sweep}) confirms that forcing a fixed hierarchy
produces \emph{half} the ablation signal of the model's preferred
allocation.  The cross-layer dimension is thus adaptive: the
architecture provides the slow field everywhere, and the optimizer
activates it where model capacity and data demand it.  The essential
scale separation is the fast/slow split \emph{within} each layer;
the cross-layer allocation is a learned resource-management decision.

The deeper question is whether attention, augmented with the right
recurrent structure, can achieve the state tracking of SSMs without
sacrificing its random-access flexibility.  Our results suggest the
answer is ``partially, and with principled guidance from the
renormalization group.''  The architecture sits between transformers
and SSMs not as a compromise but as a distinct point in the design
space---one where attention remains primary but is guided by a
multi-scale memory that respects the temporal structure of the input.


% ══════════════════════════════════════════════════════════════
\bibliographystyle{plainnat}
\begin{thebibliography}{99}

\bibitem[Boissard et~al.(2024)]{antpaper}
Boissard, E., et~al.
\newblock Pheromone trail formation at the edge of chaos: Phase diagrams of a
  lattice ant model.
\newblock 2024.

\bibitem[Dao \& Gu(2024)]{dao2024mamba2}
Dao, T. and Gu, A.
\newblock Transformers are SSMs: Generalized models and efficient algorithms
  with structured state spaces.
\newblock \emph{ICML}, 2024.

\bibitem[De et~al.(2024)]{de2024griffin}
De, S., et~al.
\newblock Griffin: Mixing gated linear recurrences with local attention for
  efficient language models.
\newblock \emph{arXiv:2402.19427}, 2024.

\bibitem[Fan et~al.(2021)]{fan2021feedback}
Fan, A., et~al.
\newblock Addressing some limitations of transformers with feedback memory.
\newblock \emph{arXiv:2002.09402}, 2021.

\bibitem[Gu et~al.(2022)]{gu2022s4}
Gu, A., Goel, K., and R\'{e}, C.
\newblock Efficiently modeling long sequences with structured state spaces.
\newblock \emph{ICLR}, 2022.

\bibitem[Gu \& Dao(2023)]{gu2023mamba}
Gu, A. and Dao, T.
\newblock Mamba: Linear-time sequence modeling with selective state spaces.
\newblock \emph{arXiv:2312.00752}, 2023.

\bibitem[Kadanoff(1966)]{kadanoff1966}
Kadanoff, L.~P.
\newblock Scaling laws for {I}sing models near $T_c$.
\newblock \emph{Physics Physique Fizika}, 2(6):263, 1966.

\bibitem[Lieber et~al.(2024)]{lieber2024jamba}
Lieber, O., et~al.
\newblock Jamba: A hybrid transformer-mamba language model.
\newblock \emph{arXiv:2403.19887}, 2024.

\bibitem[Merrill \& Sabharwal(2023)]{merrill2023transformers}
Merrill, W. and Sabharwal, A.
\newblock The parallelism tradeoff: Limitations of log-precision transformers.
\newblock \emph{TACL}, 2023.

\bibitem[Peng et~al.(2023)]{peng2023rwkv}
Peng, B., et~al.
\newblock RWKV: Reinventing RNNs for the transformer era.
\newblock \emph{EMNLP Findings}, 2023.

\bibitem[Press et~al.(2022)]{press2022alibi}
Press, O., Smith, N.~A., and Lewis, M.
\newblock Train short, test long: Attention with linear biases enables input
  length extrapolation.
\newblock \emph{ICLR}, 2022.

\bibitem[Wilson \& Kogut(1974)]{wilson1971rg}
Wilson, K.~G. and Kogut, J.
\newblock The renormalization group and the $\epsilon$ expansion.
\newblock \emph{Physics Reports}, 12(2):75--199, 1974.

\end{thebibliography}

\end{document}
